\documentclass[11pt]{article}

\usepackage[margin=1in]{geometry}
\usepackage{setspace}
\setstretch{1.05}
\usepackage{parskip}
\setlength{\parskip}{4pt}
\setlength{\parindent}{0pt}

\begin{document}

\begin{center}
\large \textbf{Reference Summary} \\
\normalsize \textbf{Shen Yan}
\end{center}

\section{chetty2014measuring}

\begin{enumerate}
  \item School districts from Washington, DC to Los Angeles have begun to calculate VA measures and use them to evaluate teachers. (Page 1)
  \item In practice, test scores from more recent classes are better predictors of current teacher quality, indicating that teacher quality fluctuates over time. (Page 2)
  \item Second, because teacher quality drifts over time, the predicted effect differs
from past performance. (Page 8)
\end{enumerate}

\section{iezzoni1997risk}

\begin{enumerate}
  \item Risks must be framed within specific time windows. (Page 4)
  \item Thus, comparisons of cost and mortality outcomes for lung cancer patients across different hospitals must account for differences in their patient mixes (Page 3)
\end{enumerate}

\section{bartholomew2011latent}

\begin{enumerate}
  \item For these reasons, any inference made from a structural model rests heavily on the substantive assumptions made about the nature of the latent variables and their relationships. (Page 236)
  \item Binary responses are extremely common, especially in the social sciences. Individuals can be classified according to whether or not they belong to a trade union or take holidays abroad. They can be recorded as agreeing or disagreeing with some proposition or as getting some item in an educational test right or wrong. (Page 95)
\end{enumerate}

\section{bond2013applying}

\begin{enumerate}
  \item a sort of within-laboratory decision to anchor the measurement scale so comparisons of person and group ability estimates taken in different research contexts might be facilitated. (Page 103)
  \item Rasch analysis software such as Winsteps allows the user to input known values (e.g., as anchored item values) so output (e.g., person values) might be more validly compared across different research contexts. (Page 104)
  \item Following on this perspective, it is clear that while Rasch folk (mostly) use DIF results in the course of ex post facto investigations and explanations of invariance failure. (Page 121)
\end{enumerate}

\section{imbens2018causal}

\begin{enumerate}
  \item In a repeated sampling thought experiment the units in the sample would have remained the same, but their assignments would might been different, leading to a different value for the estimator. (Page 4)
  \item To characterize sampling uncertainty we postulate the existence of a large, possibly infinite, population. (Page 4)
 
\end{enumerate}

\section{holland2012differential}

\begin{enumerate}
  \item Indeed, we know that individual and group differences on cognitive tests are in very large part brought about by differences in the nature, quality, and amount of education, both in school and out of school, to which people are exposed. (Page 27)
  \item Finally, the expression differential item functioning (DIF) came into use, referring to the simple observation that an item displays different statistical properties in different group settings. (Page 17)
\end{enumerate}


\end{document}


























